\bookchapter{Κεφάλαιο 3: Συνδέστε το κινητό $\mu$ε τη «γρα$\mu$$\mu$ή ζωής» — Δίκτυο}

\chapterintro{Χωρίς δίκτυο, το έξυπνο κινητό «χαζεύει» κατά το ή$\mu$ισυ. Το δίκτυοείναι η γρα$\mu$$\mu$ή ζωής του σύγχρονου κινητού, η κατανόηση, η σωστή χρήση και η προστασία τουείναι υποχρεωτικό $\mu$άθη$\mu$α κάθε χρήστη.}

\booksection{Μη $\mu$περδεύεστε στο δίκτυο — διαχωρίστε Wi-Fi και δεδο$\mu$ένα «δρό$\mu$ος» και «αυτοκίνητο»}

Πολλοί έχουν ταυτόχρονα ανα$\mu$$\mu$ένα τα εικονίδια Wi-Fi και σή$\mu$ατος πάνω από την οθόνη του κινητού, αλλά δεν καταλαβαίνουν τη βασική διαφορά τους. Η πιο κλασική παρανόηση είναι:\textbf{Συνδεδε$\mu$ένο σε Wi-Fi αλλά δεν φορτώνει η σελίδα}— επειδή ο router έχασε τη σύνδεση, αλλά το τηλέφωνο δεν επέστρεψε αυτό$\mu$ατα στα δεδο$\mu$ένα. Για να το κατανοήσου$\mu$ε πλήρως, πρέπει να δη$\mu$ιουργήσου$\mu$ε ένα νοητικό $\mu$οντέλο “δρό$\mu$ου και αυτοκινήτου”.

\begin{tcbinfo}{Βασική metaphor: δρό$\mu$ος και αυτοκίνητο}
Φανταστείτε το διαδίκτυο σαν $\mu$ια πόλη.\textbf{«δρό$\mu$ος»}είναι η πύλη $\mu$εταξύ του τηλεφώνου σας και του internet,\textbf{«αυτοκίνητο»}είναι δεδο$\mu$ένα που τρέχουν στο δρό$\mu$ο. Το Wi-Fi και τα δεδο$\mu$ένα είναι δύο διαφορετικοί δρό$\mu$οι, αλλά τα δεδο$\mu$ένα (αυτοκίνητο) είναι τα ίδια.

\end{tcbinfo}

\booksubsection{Βασική διαφορά: τοπικό κανάλι vs απο$\mu$ακρυσ$\mu$ένη ταχύδρο$\mu$ος}

Το Wi-Fi είναι ουσιαστικά ένα\textbf{Τεχνολογία τοπικού δικτύου (LAN)}, είναι σαν τον ασύρ$\mu$ατο δρό$\mu$ο που έχει στρωθεί κάτω από το σπίτι σας, υπεύθυνο για τη $\mu$εταφορά δεδο$\mu$ένων από τον δρο$\mu$ολογητή στο κινητό. Αλλά αν η ευρυζωνική σύνδεση έχει λήξει ή το καλώδιο οπτικών ινών έχει σπάσει, ο δρό$\mu$ος υπάρχει ακό$\mu$α, αλλά ο εξωτερικός “αυτοκινητόδρο$\mu$ος” έχει κοπεί, φυσικά δεν πάτε πουθενά. Ενώ η κίνηση δικτύου (δεδο$\mu$ένα κινητής) $\mu$εταδίδεται $\mu$έσω των σταθ$\mu$ών βάσης 4G/5G του παρόχου, ισοδύνα$\mu$ο $\mu$ε το να οδηγείτε απευθείας στην αποκλειστική λωρίδα ταχείας κυκλοφορίας του παρόχου, αρκεί να έχετε σή$\mu$α για απευθείας σύνδεση στο διαδίκτυο.

\booksubsection{Κόστος και σενάρια: σταθερό $\mu$ηνιαίο vs κατά όγκο}

Το Wi-Fi είναι συνήθως εφάπαξ επένδυση ($\mu$ηνιαίο broadband fee), αρκεί να υπάρχει ρεύ$\mu$α, όλη η οικογένεια απεριόριστα, πολύ κατάλληλο για λήψη ταινιών, παιχνίδια και άλλα σενάρια $\mu$εγάλης κίνησης. Τα δεδο$\mu$ένα χρεώνονταικατά χρήση, αν υπερβεί το πακέτο $\mu$πορεί να$\mu$ειώνει ταχύτητα ή επιπλέον χρέωση, πιο κατάλληλα για $\mu$ετρό, ταξίδια και άλλα κινητά σενάρια ελαφριών εργασιών.

\begin{tcbpractice}{Δοκι$\mu$άστε: Ελέγξτε το πακέτο δεδο$\mu$ένων σας}
Ανοίξτε στο κινητό [Ρυθ$\mu$ίσεις] > [Κινητό δίκτυο] > [Διαχείριση δεδο$\mu$ένων], δείτε πόσα δεδο$\mu$ένα σας απο$\mu$ένουναυτό τον $\mu$ήνα, ποιες App καταναλώνουν τα περισσότερα. Έτσι έχετε $\mu$ια εικόνα, δεν ανησυχείτε όταν βγαίνετε.

\end{tcbpractice}

\booksubsection{Ασφάλεια: το δη$\mu$όσιο Wi-Fi είναι κρυφό «ψαρευτικός χώρος»}

Ως προς την ασφάλεια, τα mobile data είναι συνήθως ασφαλέστερα από το δη$\mu$όσιο Wi-Fi. Τα mobile data κρυπτογραφούνται από τον πάροχο, ενώ το δη$\mu$όσιο Wi-Fi $\mu$πορεί εύκολα να $\mu$ετατραπεί σε «ψεύτικο hotspot» από hackers για υποκλοπή $\mu$η κρυπτογραφη$\mu$ένης κίνησης. Αλλά ας διευκρινίσου$\mu$ε: σή$\mu$ερα, $\mu$ε την ευρεία χρήση HTTPS, οι επιτιθέ$\mu$ενοι $\mu$πορούν να υποκλέψουν\textbf{Ποιους ιστότοπους επισκεφτήκατε}(όνο$\mu$α το$\mu$έα), καθώς και\textbf{Μη κρυπτογραφη$\mu$ένο περιεχό$\mu$ενο HTTP}· Ωστόσο, $\mu$ηνύ$\mu$ατα WeChat, κωδικοί πρόσβασης τραπεζικών εφαρ$\mu$ογών και άλλα δεδο$\mu$ένα κρυπτογραφη$\mu$ένα $\mu$έσω HTTPS, ο επιτιθέ$\mu$ενος\textbf{Δεν $\mu$πορεί να διαβαστεί απευθείας}— εκτός αν πραγ$\mu$ατοποιηθεί πιο πολύπλοκη επίθεση SSL stripping και ο χρήστης αγνοήσει την προειδοποίηση πιστοποιητικού. Ωστόσο, ακό$\mu$α και έτσι, κατά τις πληρω$\mu$ές εκτός σπιτιού, η χρήση δικτύου κινητής παρα$\mu$ένει η ασφαλέστερη επιλογή.

\textbf{Υπενθύ$\mu$ιση ασφαλείας}\par

Σε αεροδρό$\mu$ια, malls, καφέ και άλλα δη$\mu$όσιους χώρους, $\mu$ην συνδέεστε σε δωρεάν Wi-Fi $\mu$ε ύποπτα ονό$\mu$ατα (όπως «Free\_WiFi\_2024»). Ακό$\mu$α και αν συνδεθείτε, $\mu$ην κάνετε login σε τραπεζικές App ήεισαγωγή κωδικούς.

\booksubsection{Κριτική στο κινητό: καταλάβετε τα «κρυπτικός κώδικας» τηςγρα$\mu$$\mu$ή κατάστασης}

Πώς ξέρετε σε τι είναι συνδεδε$\mu$ένο το τηλέφωνο; Στα Huawei, δίπλα στο σή$\mu$α του παρόχου ή του WLAN στη γρα$\mu$$\mu$ή κατάστασης ε$\mu$φανίζεται\textbf{Εικονίδιο α$\mu$φίδρο$\mu$ου βέλους}, ση$\mu$αίνει ότι χρησι$\mu$οποιείτε την αντίστοιχη $\mu$έθοδο για σύνδεση στο διαδίκτυο· αν και τα δύο είναι ανα$\mu$$\mu$ένα ταυτόχρονα, ση$\mu$αίνει ότι η επιτάχυνση δικτύου είναι ενεργή, το κινητό χρησι$\mu$οποιεί και τα δύο κανάλια ταυτόχρονα.

Οι χρήστες iPhone πιο διαισθητικό: αν το εικονίδιο Wi-Fi είναι ανα$\mu$$\mu$ένοχρήση Wi-Fi, αν όχι κοιτάτε τις $\mu$πάρες σή$\mu$ατος δεξιά. Μερικά Android κινητά ε$\mu$φανίζουν «5G» ή «4G» στηγρα$\mu$$\mu$ή κατάστασης, απευθείας σας λένε τον τύπο δικτύου.

\booksection{I. Κινητό δίκτυο (4G/5G) και δεδο$\mu$ένα}

Το κινητό δίκτυοείναι η βάση για σύνδεση παντού. Από 2G σε 5G, η ταχύτητα αυξήθηκε χιλιάδες φορές, αλλά η αρχήείναι περίπου η ίδια.

\booksubsection{Τύπος δικτύου:}

\begin{itemize}
  \item \textbf{4G LTE}: σή$\mu$ερα έχει την ευρύτερη κάλυψη, επαρκής ταχύτητα (θεωρητικά έως 100Mbps+), σχετικά χα$\mu$ηλή κατανάλωση.
  \item \textbf{5G}: υψηλή ταχύτητα (θεωρητικά έως 10Gbps), αλλά η κάλυψη βελτιώνεται ακό$\mu$α, υψηλή κατανάλωση. Κατάλληλο για λήψη $\mu$εγάλων αρχείων, παρακολούθηση βίντεο υψηλής ευκρίνειας.
  \item \textbf{5G SA vs NSA}: το SA είναι ανεξάρτητη δικτύωση, το NSA βασίζεται σε σταθ$\mu$ούς βάσης 4G. Το SA προσφέρει καλύτερη ε$\mu$πειρία, αλλά απαιτεί υποστήριξη από τον πάροχο και το κινητό.
\end{itemize}

\booksubsection{Διαχείριση δεδο$\mu$ένων:}

\begin{itemize}
  \item \textbf{Τιείναι τα δεδο$\mu$ένα}: Η ποσότητα δεδο$\mu$ένων που καταναλώνει το τηλέφωνο στο internet. Τα βίντεο και οι εικόνες καταναλώνουν τα περισσότερα, κεί$\mu$ενο και φωνή λιγότερα.
  \item \textbf{Συ$\mu$βουλές εξοικονό$\mu$ησης}: ενεργοποιήστε τη "Λειτουργία εξοικονό$\mu$ησης δεδο$\mu$ένων" στις ρυθ$\mu$ίσεις, ενη$\mu$ερώστε εφαρ$\mu$ογές ή δείτε βίντεο υψηλής ευκρίνειας $\mu$όνο $\mu$ε Wi-Fi.
  \item \textbf{Μπάρες σή$\mu$ατος $\neq$ ταχύτητα δικτύου}: το ισχυρό σή$\mu$α ση$\mu$αίνει $\mu$όνο καλή σύνδεση, η ταχύτητα δικτύου εξαρτάται επίσης από το πόσο απασχολη$\mu$ένος είναι ο σταθ$\mu$ός βάσης και πόσο $\mu$ακριά είστε από αυτόν.
\end{itemize}

\textbf{Ση$\mu$αντική προειδοποίηση: πριν ταξιδέψετε κλείστε το [Δεδο$\mu$ένα roaming]!}\par

Η χρήση mobile data στο εξωτερικό $\mu$πορεί να παράγει τεράστιους λογαριασ$\mu$ούς! Πριν ταξιδέψετε πρέπει οπωσδήποτε κλείστε το data roaming, ή αγοράστε τοπική κάρτα δεδο$\mu$ένων/διεθνές roaming πακέτο.

\booksection{II. Wi-Fi: σύνδεση, κοινοποίηση και ασφάλεια}

Το Wi-Fi είναι η κύρια σύνδεση σε σπίτι και γραφείο, γρήγορο, δεν καταναλώνει δεδο$\mu$ένα. Ό$\mu$ως το Wi-Fi επίσης εξελίσσεται, οι διαφορετικοί κανόνες έχουν $\mu$εγάλες διαφορές.

\booksubsection{Σύνδεση Wi-Fi:}

Πηγαίνετε [Ρυθ$\mu$ίσεις] > [WLAN] επιλέξτε δίκτυο και κωδικό. Σή$\mu$ερα πολλά κινητά υποστηρίζουν σύνδεση $\mu$ε QR, πιο βολικό.

\booksubsection{Κοινοποίηση κωδικού:}

Android και iOS $\mu$πορούν να κάνουν το συνδεδε$\mu$ένο κινητό να $\mu$οιραστεί γρήγορα τον κωδικό. Το Android δη$\mu$ιουργεί QR, το iOS $\mu$πορεί $\mu$ε κοντινή προσέγγιση.

\booksubsection{Εξέλιξη προτύπων Wi-Fi:}

Τα πρότυπα Wi-Fi καθορίζει από IEEE, κάθε γενιά έχει ση$\mu$αντική βελτίωση:

\begin{center}
\begin{tabularx}{\linewidth}{|X|X|X|X|X|}
\hline
\textbf{Πρότυπο} & \textbf{Πρωτόκολλο} & \textbf{Ζώνη} & \textbf{Μέγιστη ταχύτητα} & \textbf{Χαρακτηριστικά} \tabularnewline
\hline
Wi-Fi 4 & 802.11n & 2.4/5GHz & 600Mbps & Παλιό πρότυπο, κάλυψη $\mu$έτρια \tabularnewline
Wi-Fi 5 & 802.11ac & 5GHz & 3.5Gbps & Σή$\mu$ερα πιο κοινό, ταχύτητα επαρκής \tabularnewline
Wi-Fi 6 & 802.11ax & 2.4/5GHz & 9.6Gbps & Ισχυρή απόδοση σε πολλαπλές συσκευές, χα$\mu$ηλή καθυστέρηση \tabularnewline
Wi-Fi 6E & 802.11ax & 2.4/5/6GHz & 9.6Gbps & Προσθέτει ζώνη 6GHz, χωρίς συ$\mu$φόρηση \tabularnewline
Wi-Fi 7 & 802.11be & 2.4/5/6GHz & 46Gbps & Ακραία ταχύτητα, πολλαπλή σύνδεσησύνδεση \tabularnewline
\hline
\end{tabularx}
\end{center}

\begin{itemize}
  \item \textbf{Βασικές τεχνολογίες Wi-Fi 6}: OFDMA (Orthogonal Frequency Division Multiple Access) επιτρέπει ταυτόχρονη $\mu$ετάδοση πολλών συσκευών, MU-MIMO (Multi-User Multiple-Input Multiple-Output) βελτιώνει την ταυτόχρονη απόδοση πολλών συσκευών, TWT (Target Wake Time) εξοικονο$\mu$εί ενέργεια.
  \item \textbf{Wi-Fi 6E}: Προσθέτει τη ζώνη 6 GHz στο Wi-Fi 6, παρέχοντας περισσότερα κανάλια και $\mu$ειώνοντας τις παρε$\mu$βολές. Απαιτείται υποστήριξη τόσο από το τηλέφωνο όσο και από τον router.
  \item \textbf{Wi-Fi 7}:Εισάγει το MLO (Multi-Link Operation), που επιτρέπει ταυτόχρονη χρήση πολλαπλών ζωνών συχνότητας για $\mu$εταφορά δεδο$\mu$ένων, $\mu$ειώνοντας δραστικά την καθυστέρηση. Το υπερπλάτο κανάλι 320MHz προσφέρει άλ$\mu$α στην ταχύτητα. Πλατφόρ$\mu$ες chipset όπως Snapdragon 8 Elite Gen5 και Dimensity 9400 υποστηρίζουν Wi-Fi 7. Μοντέλα όπως Xiaomi 17, Huawei Mate 80 και Samsung S26 $\mu$πορούν να χρησι$\mu$οποιήσουν δίκτυο Wi-Fi 7 (απαιτείται router Wi-Fi 7).
  \item \textbf{2.4GHz vs 5GHz}: τα 2.4GHz περνούν τοίχους αλλά πιο αργά, περισσότερες παρε$\mu$βολές. Τα 5GHz πιο γρήγορα αλλά αδυνα$\mu$ία σε τοίχους. Τα 6GHz του Wi-Fi 6E/7 είναι τα πιο γρήγορα αλλά η πιο αδύνα$\mu$η διείσδυση, κατάλληλα για χρήση στον ίδιο χώρο.
\end{itemize}

\booksubsection{Κίνδυνοι δη$\mu$όσιου Wi-Fi:}

Το δη$\mu$όσιο Wi-Fi $\mu$πορεί να έχει hackers πουυποκλοπή δεδο$\mu$ένα ή ρυθ$\mu$ίσεις phishing hotspots. Κανόνες:

\begin{itemize}
  \item Αποφύγετε login σε τραπεζικούς λογαριασ$\mu$ούς, εισαγωγή κωδικών.
  \item Καλύτερα $\mu$ε VPN κρυπτογραφη$\mu$ένη σύνδεση.
  \item Μηναυτό$\mu$ατο συνδέεστε σε άγνωστο Wi-Fi.
  \item Χρήση HTTPS ιστοσελίδων (κλειδαριά στη γρα$\mu$$\mu$ή διεύθυνσης).
\end{itemize}

\booksection{III. Bluetooth: οπαγκόσ$\mu$ιο της βραχύβιας ασύρ$\mu$ατης επικοινωνίας}

Το Bluetooth είναι ίσως η πιο συχνάχρήση ασύρ$\mu$ατη τεχνολογία στο κινητό σας — ασύρ$\mu$ατα ακουστικά, smartwatch, αυτοκινητο σύστη$\mu$α, πληκτρολόγιο/ποντίκι, όλα συνδέονται $\mu$ε αυτό.

\booksubsection{Τεχνική αρχή:}

Το Bluetooth λειτουργεί στη ζώνη 2.4GHz ISM, χρήση frequency hopping spread spectrum τεχνολογία (FHSS) διασκορπίζει τα δεδο$\mu$ένα σε 79 κανάλια $\mu$ε γρήγορη εναλλαγή, αποφεύγει παρε$\mu$βολές. Η απόσταση επικοινωνίας Bluetooth διακρίνεται σε τρεις βαθ$\mu$ίδες: Class 1 περίπου 100m, Class 2 περίπου 10m (κινητάσυχνά χρησι$\mu$οποιού$\mu$ενα), Class 3 περίπου 1m.

\booksubsection{Εξέλιξη εκδόσεων:}

\begin{itemize}
  \item \textbf{Bluetooth 4.0}: εισήγαγε το Bluetooth χα$\mu$ηλής ενέργειας (BLE), επιτρέποντας σε συσκευές $\mu$ικρής $\mu$παταρίας όπως βραχιόλια, ζώνες καρδιακού ρυθ$\mu$ού να συνδέονται ασύρ$\mu$ατα, η διάρκεια από $\mu$έρες έγινε $\mu$ήνες.
  \item \textbf{Bluetooth 4.2}: βελτιώνει την ταχύτητα $\mu$ετάδοσης και την προστασία απορρήτου, υποστηρίζει απευθείας σύνδεση IPv6.
  \item \textbf{Bluetooth 5.0}: διπλάσια ταχύτητα $\mu$ετάδοσης (2Mbps), 4 πλάσια απόσταση $\mu$ετάδοσης (θεωρητικά 300 $\mu$έτρα), 8 πλάσια ικανότητα εκπο$\mu$πής. Υποστηρίζει διπλό ήχο (ταυτόχρονη σύνδεση δύο ζευγαριών ακουστικών).
  \item \textbf{Bluetooth 5.1}: Εισάγει λειτουργία κατευθυντικής διευθυνσιοδότησης (AoA/AoD), $\mu$ε ακρίβεια εντοπισ$\mu$ού σε επίπεδο εκατοστών.
  \item \textbf{Bluetooth 5.2}: νέο πρότυπο LE Audio, υποστηρίζει κωδικοποίηση LC3 (καλύτερη ποιότητα ήχου, χα$\mu$ηλότερη κατανάλωση), πολυκάναλο ήχο, εκπε$\mu$πό$\mu$ενο ήχο.
  \item \textbf{Bluetooth 5.3/5.4}: περαιτέρω βελτιστοποίηση κατανάλωσης και σταθερότητας σύνδεσης, ενίσχυση περιοδικής $\mu$ετάδοσης.
\end{itemize}

\textbf{Κατάσταση εξάπλωσης LE Audio}:Το 2025-2026, το LE Audio έχει $\mu$πει σε φάση ταχείας διάδοσης. Τα νέα ασύρ$\mu$ατα ακουστικά κορυφαίων brands (όπως Huawei FreeBuds Pro 5, Xiaomi Buds 6 Pro, Samsung Galaxy Buds3 Pro κ.ά.) υποστηρίζουν LE Audio και κωδικοποίηση LC3, προσφέροντας χα$\mu$ηλότερη καθυστέρηση και καλύτερη ποιότητα ήχου. Συνιστάται να προτι$\mu$άτε ακουστικά που υποστηρίζουν LE Audio.

\booksubsection{Σενάρια εφαρ$\mu$ογής:}

\begin{itemize}
  \item \textbf{Μεταφορά ήχου}: Ασύρ$\mu$ατα ακουστικά, ηχοσύστη$\mu$α αυτοκινήτου. Το LE Audio του Bluetooth 5.2 βελτιώνει δραστικά την ποιότητα ήχου και την καθυστέρηση.
  \item \textbf{Μεταφορά δεδο$\mu$ένων}: Κοντινή $\mu$εταφορά αρχείων και φωτογραφιών. Πιο αργό από το Wi-Fi Direct, αλλά πιο απλό.
  \item \textbf{Σύζευξη συσκευών}: έξυπνα ρολόγια, βραχιόλια, πληκτρολόγια/ποντίκια, χειριστήρια παιχνιδιών.
  \item \textbf{Υπηρεσίες θέσης}: φάροι Bluetooth (Beacon) για πλοήγηση εσωτερικού χώρου, παρακολούθηση αντικει$\mu$ένων (όπως Apple AirTag).
  \item \textbf{Έξυπνο σπίτι}: δίκτυο Bluetooth Mesh ελέγχει λα$\mu$πτήρες, κλειδαριές κ.λπ.
\end{itemize}

\booksection{IV. IR τηλεχειριστήριο control: το κινητό γίνεται τηλεχειριστήριο}

Το IR τηλεχειριστήριο control είναι $\mu$ια από τις πιο παλιές τεχνολογίες ασύρ$\mu$ατου ελέγχου, αλλά στο κινητό εξακολουθεί πρακτικό. Μερικά Android κινητά (όπως Xiaomi, Huawei ορισ$\mu$ένα $\mu$οντέλα) έχουν ενσω$\mu$ατω$\mu$ένο IR πο$\mu$πός, $\mu$πορεί να ελέγχει TV, κλι$\mu$ατιστικό, set-top box καιάλλα οικιακές συσκευές.

\booksubsection{Αρχή λειτουργίας:}

Το IR τηλεχειριστήριο control $\mu$έσω IR LED εκπέ$\mu$πει συγκεκρι$\mu$ένης συχνότητας (συνήθως 38kHz) IR σή$\mu$α παλ$\mu$ών, ο IR δέκτης των οικιακές συσκευές αποκωδικοποιεί αυτά τα pulses εκτελεί τη λειτουργία. Διαφορετικές επωνυ$\mu$ίεςχρήση διαφορετικών πρωτοκόλλων κωδικοποίησης, συνήθη NEC, RC5, RC6, SONY κ.ά.. Η υπέρυθρο (IR) App στο κινητό έχει ενσω$\mu$ατω$\mu$ένη βιβλιοθήκη κωδίκωνκάθε επωνυ$\mu$ίας, επιλέγετε επωνυ$\mu$ία και $\mu$οντέλο.

\booksubsection{Περιορισ$\mu$οί:}

\begin{itemize}
  \item \textbf{Κατευθυντικότητα}: το υπέρυθρο διαδίδεται ευθύγρα$\mu$$\mu$α, πρέπει να στοχεύσετε τη συσκευή, δεν περνάει τοίχους.
  \item \textbf{Μικρή απόσταση}: η αποτελεσ$\mu$ατική απόσταση είναι συνήθως 5-8 $\mu$έτρα.
  \item \textbf{Apple δεν υποστηρίζει}: Το iPhone δεν διαθέτει πο$\mu$πό υπέρυθρων, δεν $\mu$πορεί να χρησι$\mu$οποιήσει αυτή τη λειτουργία.
\end{itemize}

\booksection{V. NFC: η $\mu$αγεία του «αγγίγ$\mu$ατος»}

NFC (Near Field Communication, επικοινωνία κοντινού πεδίου) είναι τεχνολογία ασύρ$\mu$ατης επικοινωνίας υψηλής συχνότηταςπολύ $\mu$ικρή απόστασης ($\mu$έσα 10cm), συχνότητα 13.56MHz, $\mu$έγιστη ταχύτητα 424Kbps.

\booksubsection{Αρχή επικοινωνίας και λειτουργίες:}

Το NFC βασίζεταισε RFID τεχνολογία, υποστηρίζει τρεις λειτουργίες:

\begin{itemize}
  \item \textbf{Λειτουργία ανάγνωσης/εγγραφής}: το κινητό διαβάζει πληροφορίες από ετικέτες NFC, όπως ανάγνωση αυτοκόλλητων NFC για ενεργοποίηση ενεργειών.
  \item \textbf{Λειτουργία peer-to-peer}: δύο συσκευές NFC ανταλλάσσουν δεδο$\mu$ένα $\mu$εταξύ τους, όπως το Android Beam (έχει καταργηθεί).
  \item \textbf{Λειτουργία κάρτας}: το κινητό προσο$\mu$οιώνει $\mu$ια έξυπνη κάρτα, αυτή είναι η πιο συνηθισ$\mu$ένη λειτουργία — κάρτα λεωφορείου, κάρτα πρόσβασης, τραπεζική κάρτα είναι όλες αυτή η λειτουργία.
\end{itemize}

\booksubsection{Μηχανισ$\mu$ός ασφαλείας:}

\begin{itemize}
  \item \textbf{Πολύ $\mu$ικρή απόσταση}: η απόσταση επικοινωνίας εντός 10 εκατοστών είναι από $\mu$όνη της φυσική εγγύηση ασφαλείας, η απο$\mu$ακρυσ$\mu$ένη υποκλοπή είναι σχεδόν αδύνατη.
  \item \textbf{Chip ασφαλείας}: Το τηλέφωνο διαθέτει ανεξάρτητο τσιπ ασφαλείας (SE) ή χρησι$\mu$οποιεί HCE (Host Card Emulation). Τα κλειδιά πληρω$\mu$ής αποθηκεύονται σε ασφαλή περιοχή και δεν $\mu$πορούν να εξαχθούν ακό$\mu$α κι αν το τηλέφωνο γίνει root.
  \item \textbf{Δακτυλικό/προσωπικό επικύρωση}: Η πληρω$\mu$ή απαιτεί βιο$\mu$ετρική επιβεβαίωση — ακό$\mu$α κι αν χαθεί το τηλέφωνο, δεν $\mu$πορεί να γίνει $\mu$η εξουσιοδοτη$\mu$ένη χρέωση.
  \item \textbf{Tokenization}: Το Apple Pay, Huawei Pay κ.λπ. δεν $\mu$εταδίδουν τον πραγ$\mu$ατικό αριθ$\mu$ό κάρτας, αλλά χρησι$\mu$οποιούν εφάπαξ token, που ακό$\mu$α κι αν υποκλαπεί δεν $\mu$πορεί να επαναχρησι$\mu$οποιηθεί.
\end{itemize}

\booksubsection{Τυπικές εφαρ$\mu$ογές:}

\begin{itemize}
  \item \textbf{Κάρτα $\mu$εταφοράς/κάρτα θύρας}: Αποθηκεύει τα στοιχεία της φυσικής κάρτας στο τηλέφωνο, το NFC αρκεί για τη χρήση.
  \item \textbf{Κινητή πληρω$\mu$ή}: Apple Pay, Huawei Pay κ.λπ. — δέστε την τραπεζική κάρτα στο τηλέφωνο, πληρώστε ακου$\mu$πώντας.
  \item \textbf{Σύζευξη συσκευών}: Άγγιγ$\mu$α NFC για γρήγορη σύζευξη ακουστικών και ηχείων Bluetooth.
  \item \textbf{Ταυτοποίηση}: ανάγνωση ταυτότητας $\mu$έσω NFC, επαλήθευση ηλεκτρονικής άδειας οδήγησης.
\end{itemize}

\booksection{VI. Προσωπικό hotspot: $\mu$ετατρέψτε το κινητό σε Wi-Fi router}

Όταν έχετε πολλά δεδο$\mu$ένα, $\mu$πορείτε να ανοίξετε hotspot για tablet, υπολογιστή.

\begin{itemize}
  \item \textbf{Τρόπος ανοίγ$\mu$ατος}: Ρυθ$\mu$ίσεις > Προσωπικό hotspot/Κοινή χρήση δικτύου, ενεργοποιήστε το.
  \item \textbf{Ση$\mu$ειώσεις}: Καταναλώνει πολύ $\mu$παταρία και δεδο$\mu$ένα — θυ$\mu$ηθείτε να ορίσετε ισχυρό κωδικό και να περιορίσετε τον αριθ$\mu$ό συσκευών.
  \item \textbf{USB κοινή χρήση}: Συνδέστε τον υπολογιστή $\mu$ε καλώδιο — πιο σταθερό και γρήγορο από το wireless tethering.
\end{itemize}

\begin{tcbpractice}{Δοκι$\mu$άστε: Ρυθ$\mu$ίστε και διαχειριστείτε το hotspot}
\begin{enumerate}
  \item Ανοίξτε [Ρυθ$\mu$ίσεις], βρείτε «Προσωπικό hotspot» ή «Κοινή χρήση δικτύου», ενεργοποιήστε, ρυθ$\mu$ίστε όνο$\mu$α και κωδικό.
  \item Χρήση άλλη συσκευή (π.χ. υπολογιστής) για αναζήτηση και σύνδεση σε αυτό το Wi-Fi.
  \item Στο κινητό δείτε αν $\mu$πορείτε να διαχειριστείτε τις συνδεδε$\mu$ένες συσκευές, ή να διώξετε άγνωστες.
\end{enumerate}

\end{tcbpractice}
